\section{Introduction}
Rapid urbanization and accelerating motorization have placed unprecedented strain on urban transportation systems worldwide. As cities densify, limited road infrastructure must accommodate ever-growing travel demand, leading to persistent congestion, unreliable travel times, and increased environmental externalities. Inefficient traffic flow not only wastes fuel and economic productivity but also amplifies air pollution and degrades quality of life, particularly in densely populated urban corridors. \par
Traffic signal control (TSC) aims to coordinate phase sequences and durations at signalized intersections to maximize throughput and reduce congestion. Classical fixed-time methods use the Webster formula\cite{webster} to preset signal timings based on historical traffic data collected from the traffic network at different times. In contrast, the max-pressure control framework\cite{MERCADER2020275} computes a pressure value for each signal phase from the difference between queue lengths on opposing approaches and selects the phase with the highest pressure; this strategy can be shown to maximize network throughput.\par
Reinforcement-learning (RL) methods model TSC as a Markov decision process in which an agent observes traffic states (for example, lane queue lengths and delays), chooses an action (phase switch) and receives a reward linked to throughput or delay reduction. Deep RL further parameterizes the control policy using neural networks, enabling the agent to approximate complex, high-dimensional decision functions. Such methods have demonstrated strong empirical performance across a wide range of simulated traffic scenarios. However, they also introduce notable challenges. In particular, learned policies often exhibit limited generalization, as models trained in one traffic environment exhibit degraded performance when deployed in different cities or under unseen conditions. Furthermore, the reliance on opaque neural representations reduces interpretability, making it difficult to explain or verify decision-making processes. This lack of transparency poses a barrier to real-world adoption, where reliability and accountability are critical.\par
Large Language Models (LLMs) represent a new paradigm in artificial intelligence, combining strong reasoning ability with the capacity to generate interpretable, text-based outputs. These properties make LLMs promising candidates for traffic signal control, where decision transparency and adaptability are highly desirable. By framing control as a reasoning process, LLM-based agents can potentially provide both actionable policies and human-understandable explanations, bridging the gap between performance and interpretability. However, this potential is not yet fully realized in practice. Achieving real-time responsiveness typically requires deploying smaller-scale models or mixture-of-experts (MoE) architectures with compact expert networks. Without sufficient domain adaptation, such models often fail to capture complex traffic dynamics, leading to suboptimal decisions.\par
\textbf{to answer all the demands, I ....}